\capitulo{2}{Objetivos}

El objetivo general de este Trabajo de Fin de Grado es analizar los Sistemas Aumentativos y Alternativos de Comunicación (SAAC) existentes y desarrollar una aplicación web interactiva que, mediante un cuestionario, recomiende de forma personalizada el sistema más adecuado para cada paciente.

\section{Objetivos específicos}
Para alcanzar el objetivo general, se plantean los siguientes objetivos específicos:
\begin{itemize}
    	\item \textbf{Analizar el catálogo tecnológico:} Estudiar las soluciones de SAAC disponibles en el mercado, indexando sus especificaciones de hardware, software y métodos de acceso.
    	\item \textbf{Definir las variables clínicas:} Identificar los parámetros de capacidad del usuario (motora, cognitiva, visual y lingüística) que determinan la elección de un SAAC.
    	\item \textbf{Diseñar el modelo de recomendación:} Desarrollar la lógica que vincula las respuestas del cuestionario del usuario con las características de los sistemas del catálogo.
    	\item \textbf{Desarrollar un prototipo funcional:} Construir la aplicación web interactiva que implemente el cuestionario, procese los filtros y recoja feedback de forma anónima.
    	\item \textbf{Evaluar limitaciones del sistema:} Analizar de forma crítica los puntos débiles del prototipo, su infraestructura y su viabilidad legal de cara a un uso clínico real.
\end{itemize}

\section{Objetivos técnicos}
Desde la perspectiva de la ingeniería del software, se establecen los siguientes objetivos técnicos:
\begin{itemize}
    	\item \textbf{Implementar la base de datos:} Modelar y construir una base de datos relacional en PostgreSQL utilizando el ORM SQLAlchemy para almacenar el catálogo de SAAC.
    	\item \textbf{Diseñar la arquitectura del software:} Estructurar la aplicación web utilizando el framework Flask y el motor de plantillas Jinja2 para la interfaz de usuario.
    	\item \textbf{Programar el motor de reglas:} Implementar la lógica de filtrado por descarte mediante condicionales IF-THEN en Python, incluyendo el soporte para Voice Banking.
    	\item \textbf{Integrar medidas de seguridad web:} Aplicar cabeceras de seguridad y mitigar vulnerabilidades comunes utilizando la librería Flask-Talisman bajo estándares OWASP.
    	\item \textbf{Optimizar la interfaz y accesibilidad:} Diseñar las pantallas (HTML / CSS) para facilitar la navegación a usuarios no especializados y asentar las bases para el uso de pulsadores o eyetracking.
\end{itemize}

\section{Objetivos formativos}
Finalmente, desde el punto de vista del aprendizaje, el trabajo persigue los siguientes objetivos:
\begin{itemize}
   	 \item \textbf{Aplicar la Ingeniería de la Salud:} Resolver un problema real del ámbito sanitario combinando el desarrollo de software con criterios médicos.
    	\item \textbf{Gestionar el ciclo de vida del proyecto:} Adquirir competencias en análisis de requisitos, presupuestos bajo el Convenio TIC y cumplimiento legal del RGPD.
    	\item \textbf{Desarrollar autocrítica técnica:} Aprender a identificar la deuda técnica del proyecto y proponer soluciones de mejora.
\end{itemize}